\section{Рабочий проект}

Разработанная система представляет собой графическое приложение для настройки клавиатур с поддержкой расширенных протоколов динамической конфигурации. Ниже приведено описание основных компонентов (модулей) системы, их файловая структура и взаимосвязи.

\subsection{Пакет \texttt{autodetect}}

\subsubsection{Модуль \texttt{device\_monitor.py}}

Реализует класс \texttt{DeviceMonitor} (синглтон), который управляет списком доступных HID-устройств, поддерживающих расширенный протокол, и обеспечивает переключение между ними. Использует отдельный поток \texttt{MonitorThread} для фонового сканирования.

Методы класса \texttt{DeviceMonitor} представлены в таблице \ref{tab:devicemonitor}.

\begin{xltabular}{\textwidth}{|p{3.5cm}|p{3.5cm}|>{\setlength{\baselineskip}{0.7\baselineskip}}X|}
	\caption{Методы класса \texttt{DeviceMonitor}\label{tab:devicemonitor}}\\
	\hline \centrow Название метода & \centrow Аргументы & \centrow Описание \\
	\hline
	\endfirsthead
	\caption*{Продолжение таблицы \ref{tab:devicemonitor}}\\
	\hline \centrow Название метода & \centrow Аргументы & \centrow Описание \\
	\hline
	\finishhead
	lock & Нет аргументов & Возвращает объект \texttt{MonitorLocker} для исключающей блокировки операций обновления. \\
	\hline
	load\_dummy & data & Загружает стороннее JSON-определение клавиатуры (dummy-режим). \\
	\hline
	sideload\_legacy\_\newline json & data & Загружает JSON для legacy-клавиатуры. \\
	\hline
	load\_legacy\_\newline stack & data & Загружает стек legacy-определений (VIA). \\
	\hline
	select\_device & idx & Выбирает устройство по индексу, открывает соединение и перезагружает конфигурацию. \\
	\hline
	on\_devices\_\newline updated & devices, changed & Слот, обновляющий список устройств в главном окне. \\
	\hline
	update & quiet, hard & Запускает сканирование устройств (вызов \texttt{MonitorThread.update}). \\
	\hline
\end{xltabular}

\subsubsection{Модуль \texttt{monitor\_thread.py}}

Содержит класс \texttt{MonitorThread}, производный от \texttt{QThread}. Выполняет периодическое (раз в секунду) сканирование HID-устройств через \texttt{find\_custom\_devices}. Сигнал \texttt{devices\_updated} передаёт обновлённый список в главный поток.

Методы класса \texttt{MonitorThread} представлены в таблице \ref{tab:monitorthread}.

\begin{xltabular}{\textwidth}{|p{3.5cm}|p{3.5cm}|>{\setlength{\baselineskip}{0.7\baselineskip}}X|}
	\caption{Методы класса \texttt{MonitorThread}\label{tab:monitorthread}}\\
	\hline \centrow Название метода & \centrow Аргументы & \centrow Описание \\
	\hline
	\endfirsthead
	\caption*{Продолжение таблицы \ref{tab:monitorthread}}\\
	\hline \centrow Название метода & \centrow Аргументы & \centrow Описание \\
	\hline
	\finishhead
	run & Нет аргументов & Основной цикл потока: вызов \texttt{update}, пауза 1 с. \\
	\hline
	lock / unlock & Нет аргументов & Блокировка/разблокировка обновлений (используется при длительных операциях). \\
	\hline
	update & quiet, hard & Сканирует устройства (вне мьютекса), затем обновляет внутренний список. \\
	\hline
	load\_dummy / sideload\_legacy\_\newline json / load\_legacy\_\newline stack & data & Сохраняют переданные JSON-данные для дальнейшего использования. \\
	\hline
	set\_device & current\_device & Устанавливает текущее выбранное устройство. \\
	\hline
\end{xltabular}

\subsubsection{Модуль \texttt{device\_handler.py}}

Определяет классы для работы с различными типами устройств: \texttt{BaseDevice} (общий), \texttt{CustomKeyboard} (основная клавиатура), \texttt{BootloaderDevice} (режим загрузчика), \texttt{DummyKeyboard} (эмуляция). Содержит логику открытия/закрытия HID-канала и получения UID.

Методы классов \texttt{BaseDevice}, \texttt{CustomKeyboard}, \texttt{BootloaderDevice} представлены в таблице \ref{tab:devicehandler}.

\begin{xltabular}{\textwidth}{|p{3.5cm}|p{3.5cm}|>{\setlength{\baselineskip}{0.7\baselineskip}}X|}
	\caption{Методы классов \texttt{BaseDevice}, \texttt{CustomKeyboard}, \texttt{BootloaderDevice}\label{tab:devicehandler}}\\
	\hline \centrow Название метода & \centrow Аргументы & \centrow Описание \\
	\hline
	\endfirsthead
	\caption*{Продолжение таблицы \ref{tab:devicehandler}}\\
	\hline \centrow Название метода & \centrow Аргументы & \centrow Описание \\
	\hline
	\finishhead
	open & override\_json & Открывает устройство по пути; для \texttt{CustomKeyboard} загружает определение. \\
	\hline
	send / recv & data, length & Низкоуровневый обмен HID-отчётами. \\
	\hline
	close & Нет аргументов & Закрывает HID-устройство. \\
	\hline
	title & Нет аргументов & Возвращает строку для отображения в комбобоксе. \\
	\hline
	get\_uid & Нет аргументов & Запрашивает UID клавиатуры (или загрузчика). \\
	\hline
\end{xltabular}

\subsubsection{Модуль \texttt{hid\_comm.py}}

Предоставляет унифицированный интерфейс к HID-библиотеке (на Linux используется \texttt{hidraw}, на остальных – стандартный \texttt{hid}). Содержит функцию \texttt{hid\_send}, которая добавляет нулевой префикс (report ID) и обрабатывает тайм-ауты.

Функции модуля \texttt{hid\_comm.py} представлены в таблице \ref{hid-comm:table}.

\begin{xltabular}{\textwidth}{|p{3.5cm}|p{3.5cm}|>{\setlength{\baselineskip}{0.7\baselineskip}}X|}
	\caption{Функции модуля \texttt{hid\_comm.py}\label{hid-comm:table}}\\
	\hline \centrow Название функции & \centrow Аргументы & \centrow Описание \\
	\hline
	\finishhead
	hid\_send & dev, msg, retries & Отправляет сообщение длиной до 32 байт; при неудаче повторяет попытки. \\
	\hline
\end{xltabular}

\subsection{Пакет \texttt{comm}}

\subsubsection{Модуль \texttt{base\_comm.py}}

Содержит абстрактный класс \texttt{BaseProtocol} с методом \texttt{\_retrieve\_dynamic\_entries} для пакетного чтения динамических данных (tap dance, combos и т.п.) от клавиатуры.

Метод класса \texttt{BaseProtocol} представлены в таблице \ref{BaseProtocol:table}.

\begin{xltabular}{\textwidth}{|p{3.5cm}|p{3.5cm}|>{\setlength{\baselineskip}{0.7\baselineskip}}X|}
	\caption{Метод класса \texttt{BaseProtocol}\label{BaseProtocol:table}}\\
	\hline \centrow Название метода & \centrow Аргументы & \centrow Описание \\
	\hline
	\finishhead
	\_retrieve\_\newline dynamic\_\newline entries & cmd, count, fmt & Последовательно запрашивает \texttt{count} записей, распаковывает их согласно формату \texttt{fmt}. \\
	\hline
\end{xltabular}

\subsubsection{Модуль \texttt{constants.py}}

Хранит числовые константы протоколов: команды \texttt{CMD\_*}, коды динамических операций \texttt{DYNAMIC\_*}, версии протоколов \texttt{PROTOCOL\_*} и размер буфера \texttt{BUFFER\_FETCH\_CHUNK}. Функций не содержит.

\subsubsection{Модуль \texttt{keyboard\_comm.py}}

Главный класс \texttt{Keyboard}, объединяющий все протоколы (через множественное наследование). Обеспечивает загрузку определения клавиатуры (JSON + LZMA), чтение/запись раскладки, энкодеров, настроек подсветки, а также команды разблокировки/блокировки.

Методы класса \texttt{Keyboard} представлены в таблице \ref{tab:keyboard}.

\begin{xltabular}{\textwidth}{|p{3.5cm}|p{3.5cm}|>{\setlength{\baselineskip}{0.7\baselineskip}}X|}
	\caption{Методы класса \texttt{Keyboard}\label{tab:keyboard}}\\
	\hline \centrow Название метода & \centrow Аргументы & \centrow Описание \\
	\hline
	\endfirsthead
	\caption*{Продолжение таблицы \ref{tab:keyboard}}\\
	\hline \centrow Название метода & \centrow Аргументы & \centrow Описание \\
	\hline
	\finishhead
	reload & sideload\_json & Загружает все параметры клавиатуры: слои, матрицу, раскладку, макросы, tap dance, комбинации, переопределения. \\
	\hline
	set\_key & layer, row, col, code & Записывает код клавиши в указанный слой и позицию матрицы. \\
	\hline
	set\_encoder & layer, index, direction, code & Устанавливает код для энкодера (0 – CCW, 1 – CW). \\
	\hline
	save\_layout & Нет аргументов & Сериализует раскладку, макросы, tap dance, комбинации, настройки QMK в JSON-строку. \\
	\hline
	restore\_layout & data & Восстанавливает ранее сохранённую раскладку из JSON. \\
	\hline
	reset & Нет аргументов & Отправляет команду перезагрузки клавиатуры в режим загрузчика. \\
	\hline
	get\_unlock\_status & retries & Возвращает статус разблокировки. \\
	\hline
	unlock\_start / unlock\_poll & Нет аргументов & Инициирует и опрашивает процесс разблокировки. \\
	\hline
	lock & Нет аргументов & Блокирует клавиатуру (запрещает динамическую запись). \\
	\hline
	matrix\_poll & Нет аргументов & Запрашивает состояние матрицы (для тестера). \\
	\hline
	qmk\_settings\_set & qsid, value & Устанавливает значение настройки QMK. \\
	\hline
\end{xltabular}

\subsubsection{Модуль \texttt{macro\_protocol.py}}

Реализует сериализацию/десериализацию макросов в бинарный формат (версии 1 и 2). Содержит методы \texttt{macro\_serialize}, \texttt{macro\_deserialize}, \texttt{set\_macro}, \texttt{macros\_serialize}, \texttt{macros\_deserialize}. Поддерживает действия \texttt{ActionText}, \texttt{ActionTap}, \texttt{ActionDown}, \texttt{ActionUp}, \texttt{ActionDelay}.

Методы протокола макросов представлены в таблице \ref{macro-protocol:table}.

\begin{xltabular}{\textwidth}{|p{3.5cm}|p{3.5cm}|>{\setlength{\baselineskip}{0.7\baselineskip}}X|}
	\caption{Методы протокола макросов\label{macro-protocol:table}}\\
	\hline \centrow Название метода & \centrow Аргументы & \centrow Описание \\
	\hline
	\finishhead
	reload\_macros\_\newline early & Нет аргументов & Запрашивает количество макросов и размер буфера. \\
	\hline
	reload\_macros\_\newline late & Нет аргументов & Загружает бинарный буфер макросов с устройства. \\
	\hline
	set\_macro & data & Отправляет сериализованные макросы на устройство (с разблокировкой). \\
	\hline
	macro\_serialize / macro\_deserialize & macro / data & Преобразование отдельного макроса. \\
	\hline
	macros\_serialize / macros\_\newline deserialize & macros / data & Преобразование списка всех макросов. \\
	\hline
\end{xltabular}

\subsubsection{Модуль \texttt{tap\_dance\_protocol.py}}

Предоставляет методы \texttt{reload\_tap\_dance}, \texttt{tap\_dance\_get}, \texttt{tap\_dance\_set}, \texttt{save\_tap\_dance}, \texttt{restore\_tap\_dance}. Работает с динамическими командами \texttt{DYNAMIC\_TAP\_DANCE\_GET/SET}. Каждая запись содержит четыре действия (tap, hold, double-tap, tap+hold) и время тап-термина.

Методы протокола Tap Dance представлены в таблице \ref{tap-dance-protocol:table}.

\begin{xltabular}{\textwidth}{|p{3.5cm}|p{3.5cm}|>{\setlength{\baselineskip}{0.7\baselineskip}}X|}
	\caption{Методы протокола Tap Dance\label{tap-dance-protocol:table}}\\
	\hline \centrow Название метода & \centrow Аргументы & \centrow Описание \\
	\hline
	\finishhead
	reload\_tap\_dance & Нет аргументов & Загружает все tap dance записи с устройства. \\
	\hline
	tap\_dance\_get & idx & Возвращает запись по индексу. \\
	\hline
	tap\_dance\_set & idx, entry & Отправляет изменённую запись на устройство. \\
	\hline
	save\_tap\_dance & Нет аргументов & Сохраняет все записи в список для JSON. \\
	\hline
	restore\_tap\_dance & data & Восстанавливает записи из JSON. \\
	\hline
\end{xltabular}

\subsubsection{Модуль \texttt{combo\_protocol.py}}

Аналогично предыдущему, но для комбинаций клавиш (до четырёх триггеров и выходной код). Использует \texttt{DYNAMIC\_COMBO\_GET/SET}.

Методы протокола комбинаций представлены в таблице \ref{combo-protocol:table}.

\begin{xltabular}{\textwidth}{|p{3.5cm}|p{3.5cm}|>{\setlength{\baselineskip}{0.7\baselineskip}}X|}
	\caption{Методы протокола комбинаций\label{combo-protocol:table}}\\
	\hline \centrow Название метода & \centrow Аргументы & \centrow Описание \\
	\hline
	\finishhead
	reload\_combo & Нет аргументов & Загружает все комбинации с устройства. \\
	\hline
	combo\_get & idx & Возвращает комбинацию по индексу. \\
	\hline
	combo\_set & idx, entry & Отправляет изменённую комбинацию на устройство. \\
	\hline
	save\_combo & Нет аргументов & Сохраняет все комбинации в список для JSON. \\
	\hline
	restore\_combo & data & Восстанавливает комбинации из JSON. \\
	\hline
\end{xltabular}

\subsubsection{Модуль \texttt{key\_override\_protocol.py}}

Управляет переопределениями клавиш (key overrides). Каждая запись содержит: триггерную клавишу, замену, маски модификаторов (требуемые, отрицательные, подавляемые), битовую маску слоёв и опции.

Методы протокола переопределений клавиш представлены в таблице \ref{key-override-protocol:table}.

\begin{xltabular}{\textwidth}{|p{3.5cm}|p{3.5cm}|>{\setlength{\baselineskip}{0.7\baselineskip}}X|}
	\caption{Методы протокола переопределений клавиш\label{key-override-protocol:table}}\\
	\hline \centrow Название метода & \centrow Аргументы & \centrow Описание \\
	\hline
	\finishhead
	reload\_key\_\newline override & Нет аргументов & Загружает все переопределения с устройства. \\
	\hline
	key\_override\_get & idx & Возвращает запись по индексу. \\
	\hline
	key\_override\_set & idx, entry & Отправляет изменённую запись на устройство. \\
	\hline
	save\_key\_\newline override & Нет аргументов & Сохраняет все записи в список для JSON. \\
	\hline
	restore\_key\_\newline override & data & Восстанавливает записи из JSON. \\
	\hline
\end{xltabular}

\subsubsection{Модуль \texttt{alt\_repeat\_key.py}}

Реализует протокол альтернативного повтора (Alt Repeat Key). Запись содержит: клавишу для повтора, альтернативную клавишу, маску разрешённых модификаторов и опции.

Методы протокола альтернативного повтора представлены в таблице \ref{alt-repeat-key:table}.

\begin{xltabular}{\textwidth}{|p{3.5cm}|p{3.5cm}|>{\setlength{\baselineskip}{0.7\baselineskip}}X|}
	\caption{Методы протокола альтернативного повтора\label{alt-repeat-key:table}}\\
	\hline \centrow Название метода & \centrow Аргументы & \centrow Описание \\
	\hline
	\finishhead
	reload\_alt\_repeat & Нет аргументов & Загружает все записи с устройства. \\
	\hline
	alt\_repeat\_get & idx & Возвращает запись по индексу. \\
	\hline
	alt\_repeat\_set & idx, entry & Отправляет изменённую запись на устройство. \\
	\hline
	save\_alt\_repeat & Нет аргументов & Сохраняет все записи в список для JSON. \\
	\hline
	restore\_alt\_repeat & data & Восстанавливает записи из JSON. \\
	\hline
\end{xltabular}

\subsubsection{Модуль \texttt{dynamic\_protocol.py}}

Отвечает за получение списка поддерживаемых динамических возможностей: \texttt{reload\_dynamic} запрашивает у устройства количество tap dance, комбинаций, key overrides и alt repeat entries, а также битовую маску дополнительных фич (\texttt{caps\_word}, \texttt{layer\_lock} и т.п.).

Метод протокола динамических возможностей представлены в таблице \ref{dynamic-protocol:table}.

\begin{xltabular}{\textwidth}{|p{3.5cm}|p{3.5cm}|>{\setlength{\baselineskip}{0.7\baselineskip}}X|}
	\caption{Метод протокола динамических возможностей\label{dynamic-protocol:table}}\\
	\hline \centrow Название метода & \centrow Аргументы & \centrow Описание \\
	\hline
	\finishhead
	reload\_dynamic & Нет аргументов & Запрашивает у устройства количество записей каждого типа и поддерживаемые фичи. \\
	\hline
\end{xltabular}

\subsection{Пакет \texttt{config}}

\subsubsection{Модуль \texttt{basic\_config.py}}

Абстрактный базовый класс \texttt{BasicConfig} (наследник \texttt{QVBoxLayout}) для всех редакторов. Определяет методы \texttt{valid()}, \texttt{rebuild(device)}, \texttt{activate()}, \texttt{deactivate()}.

Методы класса \texttt{BasicConfig} представлены в таблице \ref{BasicConfig:table}.

\begin{xltabular}{\textwidth}{|p{3.5cm}|p{3.5cm}|>{\setlength{\baselineskip}{0.7\baselineskip}}X|}
	\caption{Методы класса \texttt{BasicConfig}\label{BasicConfig:table}}\\
	\hline \centrow Название метода & \centrow Аргументы & \centrow Описание \\
	\hline
	\finishhead
	valid & Нет аргументов & Возвращает \texttt{True}, если редактор может работать с текущим устройством. \\
	\hline
	rebuild & device & Перестраивает интерфейс на основе данных устройства. \\
	\hline
	activate / deactivate & Нет аргументов & Вызывается при переключении вкладки для захвата/освобождения ресурсов. \\
	\hline
\end{xltabular}

\subsubsection{Модуль \texttt{keymap\_editor.py}}

Класс \texttt{KeymapEditor} отображает графическое представление клавиатуры (\texttt{KeyboardWidget}), позволяет переключать слои, выбирать отдельные клавиши и назначать им коды.

Методы класса \texttt{KeymapEditor} представлены в таблице \ref{tab:keymapeditor}.

\begin{xltabular}{\textwidth}{|p{3.5cm}|p{3.5cm}|>{\setlength{\baselineskip}{0.7\baselineskip}}X|}
	\caption{Методы класса \texttt{KeymapEditor}\label{tab:keymapeditor}}\\
	\hline \centrow Название метода & \centrow Аргументы & \centrow Описание \\
	\hline
	\endfirsthead
	\caption*{Продолжение таблицы \ref{tab:keymapeditor}}\\
	\hline \centrow Название метода & \centrow Аргументы & \centrow Описание \\
	\hline
	\finishhead
	rebuild & device & Создаёт кнопки слоёв, передаёт данные о клавишах в \texttt{KeyboardWidget}. \\
	\hline
	set\_key & keycode & Присваивает выбранной клавише указанный код (поддерживает маскированные коды). \\
	\hline
	switch\_layer & idx & Переключает текущий отображаемый слой. \\
	\hline
	refresh\_layer\_\newline display & Нет аргументов & Обновляет текст на всех клавишах согласно текущему слою и переопределениям раскладки. \\
	\hline
	on\_keycode\_\newline changed & code & Обработчик выбора кода из панели \texttt{KeycodeTabs}. \\
	\hline
	save\_layout & Нет аргументов & Возвращает сериализованную раскладку (вызов \texttt{Keyboard.save\_layout}). \\
	\hline
	restore\_layout & data & Восстанавливает раскладку из JSON. \\
	\hline
\end{xltabular}

\subsubsection{Модуль \texttt{layout\_editor.py}}

Класс \texttt{LayoutEditor} управляет аппаратными опциями раскладки (layout options), которые задаются в JSON-определении (флаги типа boolean или выбор из списка). Опции сохраняются в \texttt{Keyboard.layout\_options} и влияют на отображение переключаемых частей клавиатуры.

Методы класса \texttt{LayoutEditor} представлены в таблице \ref{LayoutEditor:table}.

\begin{xltabular}{\textwidth}{|p{3.5cm}|p{3.5cm}|>{\setlength{\baselineskip}{0.7\baselineskip}}X|}
	\caption{Методы класса \texttt{LayoutEditor}\label{LayoutEditor:table}}\\
	\hline \centrow Название метода & \centrow Аргументы & \centrow Описание \\
	\hline
	\finishhead
	rebuild & device & Создаёт элементы управления (чекбоксы, комбобоксы) на основе \texttt{layout\_labels}. \\
	\hline
	pack & Нет аргументов & Возвращает упакованные в целое число значения всех опций. \\
	\hline
	unpack & value & Устанавливает опции из целого числа. \\
	\hline
	get\_choice & index & Возвращает значение опции по индексу. \\
	\hline
\end{xltabular}

\subsubsection{Модуль \texttt{macro\_recorder.py}}

Класс \texttt{MacroRecorder} (вкладка «Macros») позволяет создавать и редактировать макросы, записывать их с клавиатуры, оптимизировать. Использует виджеты \texttt{MacroTab} и \texttt{MacroLine}.

Методы класса \texttt{MacroRecorder} представлены в таблице \ref{MacroRecorder:table}.

\begin{xltabular}{\textwidth}{|p{3.5cm}|p{3.5cm}|>{\setlength{\baselineskip}{0.7\baselineskip}}X|}
	\caption{Методы класса \texttt{MacroRecorder}\label{MacroRecorder:table}}\\
	\hline \centrow Название метода & \centrow Аргументы & \centrow Описание \\
	\hline
	\finishhead
	rebuild & device & Создаёт вкладки для каждого макроса, загружает текущие макросы. \\
	\hline
	serialize / deserialize & Нет аргументов / data & Преобразование макросов в бинарный буфер и обратно. \\
	\hline
	on\_record & tab, append & Начинает интерактивную запись макроса. \\
	\hline
	on\_keystroke & keystroke & Преобразует перехваченное нажатие в \texttt{KeyDown/Up/Tap/String}. \\
	\hline
	on\_save & Нет аргументов & Разблокирует клавиатуру и отправляет макросы. \\
	\hline
\end{xltabular}

\subsubsection{Модуль \texttt{combo\_editor.py}}

Класс \texttt{ComboEditor} (вкладка «Combos») позволяет настраивать комбинации из нескольких одновременно нажатых клавиш.

Методы класса \texttt{ComboEditor} представлены в таблице \ref{ComboEditor:table}.

\begin{xltabular}{\textwidth}{|p{3.5cm}|p{3.5cm}|>{\setlength{\baselineskip}{0.7\baselineskip}}X|}
	\caption{Методы класса \texttt{ComboEditor}\label{ComboEditor:table}}\\
	\hline \centrow Название метода & \centrow Аргументы & \centrow Описание \\
	\hline
	\finishhead
	rebuild\_ui & Нет аргументов & Создаёт вкладки для каждой поддерживаемой комбинации. \\
	\hline
	load & data & Загружает комбинацию из кортежа (4 клавиши-триггера + выходная клавиша). \\
	\hline
	save & Нет аргументов & Возвращает текущие значения комбинации в виде кортежа. \\
	\hline
	on\_key\_changed & Нет аргументов & Отправляет изменённую комбинацию на устройство через \texttt{Keyboard.combo\_set}. \\
	\hline
\end{xltabular}

\subsubsection{Модуль \texttt{tap\_dance\_editor.py}}

Класс \texttt{TapDanceEditor} (вкладка «Tap Dance») управляет действиями на одно нажатие, двойное нажатие, удержание и tap+hold.

Методы класса \texttt{TapDanceEditor} представлены в таблице \ref{TapDanceEditor:table}.

\begin{xltabular}{\textwidth}{|p{3.5cm}|p{3.5cm}|>{\setlength{\baselineskip}{0.7\baselineskip}}X|}
	\caption{Методы класса \texttt{TapDanceEditor}\label{TapDanceEditor:table}}\\
	\hline \centrow Название метода & \centrow Аргументы & \centrow Описание \\
	\hline
	\finishhead
	rebuild\_ui & Нет аргументов & Создаёт вкладки для каждого tap dance. \\
	\hline
	load & data & Загружает данные степа: четыре действия и значение тайминга. \\
	\hline
	save & Нет аргументов & Возвращает кортеж (действие на тап, на удержание, на двойной тап, на тап+удержание, тайминг). \\
	\hline
	on\_save & Нет аргументов & Отправляет все изменённые tap dance на устройство. \\
	\hline
	on\_revert & Нет аргументов & Отменяет локальные изменения, перезагружая данные с клавиатуры. \\
	\hline
\end{xltabular}

\subsubsection{Модуль \texttt{key\_override\_editor.py}}

Класс \texttt{KeyOverride} (вкладка «Key Overrides») настраивает условную замену клавиш с учётом модификаторов и слоёв.

Методы класса \texttt{KeyOverride} представлены в таблице \ref{KeyOverride:table}.

\begin{xltabular}{\textwidth}{|p{3.5cm}|p{3.5cm}|>{\setlength{\baselineskip}{0.7\baselineskip}}X|}
	\caption{Методы класса \texttt{KeyOverride}\label{KeyOverride:table}}\\
	\hline \centrow Название метода & \centrow Аргументы & \centrow Описание \\
	\hline
	\finishhead
	rebuild\_ui & Нет аргументов & Создаёт вкладки для каждого переопределения (до 128). \\
	\hline
	load & ko & Загружает запись переопределения. \\
	\hline
	save & Нет аргументов & Возвращает объект \texttt{KeyOverrideEntry}. \\
	\hline
	on\_change & Нет аргументов & Отправляет изменённое переопределение на устройство. \\
	\hline
\end{xltabular}

\subsubsection{Модуль \texttt{alt\_key\_repeat.py}}

Класс \texttt{AltRepeatKey} (вкладка «Alt Repeat Key») конфигурирует альтернативный повтор: клавиша, заменяющая последнюю нажатую, разрешённые модификаторы и опции.

Методы класса \texttt{AltRepeatKey} представлены в таблице \ref{AltRepeatKey:table}.

\begin{xltabular}{\textwidth}{|p{3.5cm}|p{3.5cm}|>{\setlength{\baselineskip}{0.7\baselineskip}}X|}
	\caption{Методы класса \texttt{AltRepeatKey}\label{AltRepeatKey:table}}\\
	\hline \centrow Название метода & \centrow Аргументы & \centrow Описание \\
	\hline
	\finishhead
	rebuild\_ui & Нет аргументов & Создаёт вкладки для каждой альтернативной клавиши повтора. \\
	\hline
	load & arep & Загружает запись альтернативного повтора. \\
	\hline
	save & Нет аргументов & Возвращает объект \texttt{AltRepeatEntry}. \\
	\hline
	on\_change & Нет аргументов & Отправляет изменения на устройство. \\
	\hline
\end{xltabular}

\subsubsection{Модуль \texttt{rgb\_control.py}}

Класс \texttt{RGBControl} объединяет управление тремя типами подсветки: QMK Backlight, QMK RGBLight и Custom RGB.

Методы класса \texttt{RGBControl} представлены в таблице \ref{RGBControl:table}.

\begin{xltabular}{\textwidth}{|p{3.5cm}|p{3.5cm}|>{\setlength{\baselineskip}{0.7\baselineskip}}X|}
	\caption{Методы класса \texttt{RGBControl}\label{RGBControl:table}}\\
	\hline \centrow Название метода & \centrow Аргументы & \centrow Описание \\
	\hline
	\finishhead
	rebuild & device & Определяет поддерживаемые типы подсветки и создаёт обработчики. \\
	\hline
	update\_from\_\newline keyboard & Нет аргументов & Запрашивает текущие значения с устройства и обновляет интерфейс. \\
	\hline
	set\_qmk\_rgblight\newline \_brightness / effect & value, index & Отправляет команды изменения яркости/эффекта RGBLight. \\
	\hline
	set\_custom\_rgb\_\newline mode / speed / color & value & Управляет Custom RGB. \\
	\hline
	save\_rgb & Нет аргументов & Сохраняет настройки подсветки в EEPROM клавиатуры. \\
	\hline
\end{xltabular}

\subsubsection{Модуль \texttt{qmk\_settings\_editor.py}}

Класс \texttt{QmkSettingsEditor} динамически строит интерфейс на основе файла \texttt{qmk\_settings.json}. Поддерживает boolean и integer настройки.

Методы класса \texttt{QmkSettingsEditor} представлены в таблице \ref{qmkseteditor:table}.

\begin{xltabular}{\textwidth}{|p{3.5cm}|p{3.5cm}|>{\setlength{\baselineskip}{0.7\baselineskip}}X|}
	\caption{Методы класса \texttt{QmkSettingsEditor}\label{qmkseteditor:table}}\\
	\hline \centrow Название метода & \centrow Аргументы & \centrow Описание \\
	\hline
	\finishhead
	recreate\_gui & Нет аргументов & Создаёт вкладки и элементы управления согласно \texttt{supported\_settings}. \\
	\hline
	reload\_settings & Нет аргументов & Запрашивает все настройки с клавиатуры и обновляет интерфейс. \\
	\hline
	prepare\_settings & Нет аргументов & Собирает словарь QSID → значение из текущих виджетов. \\
	\hline
	save\_settings & Нет аргументов & Отправляет изменённые настройки на устройство. \\
	\hline
	reset\_settings & Нет аргументов & Сбрасывает все настройки QMK к значениям по умолчанию. \\
	\hline
\end{xltabular}

\subsubsection{Модуль \texttt{matrix\_tester.py}}

Класс \texttt{MatrixTester} (вкладка «Matrix tester») периодически опрашивает состояние матрицы клавиш и отображает нажатые клавиши.

Методы класса \texttt{MatrixTester} представлены в таблице \ref{matrixtester:table}.

\begin{xltabular}{\textwidth}{|p{3.5cm}|p{3.5cm}|>{\setlength{\baselineskip}{0.7\baselineskip}}X|}
	\caption{Методы класса \texttt{MatrixTester}\label{matrixtester:table}}\\
	\hline \centrow Название метода & \centrow Аргументы & \centrow Описание \\
	\hline
	\finishhead
	rebuild & device & Загружает матрицу из определения клавиатуры. \\
	\hline
	matrix\_poller & Нет аргументов & Вызывается по таймеру, получает состояние матрицы и обновляет виджеты. \\
	\hline
	unlock & Нет аргументов & Вызывает диалог разблокировки. \\
	\hline
	activate / deactivate & Нет аргументов & Захватывает/освобождает ввод и управляет таймером. \\
	\hline
\end{xltabular}

\subsubsection{Модуль \texttt{firmware\_updater.py}}

Реализует прошивку клавиатуры из файла \texttt{.vfw}. Проверяет сигнатуру, SHA256, UID, отправляет данные блоками по 64 байта в загрузчик.

Функции модуля \texttt{firmware\_updater.py} представлены в таблице \ref{firmware-updater:table}.

\begin{xltabular}{\textwidth}{|p{3.5cm}|p{3.5cm}|>{\setlength{\baselineskip}{0.7\baselineskip}}X|}
	\caption{Функции модуля \texttt{firmware\_updater.py}\label{firmware-updater:table}}\\
	\hline \centrow Название функции & \centrow Аргументы & \centrow Описание \\
	\hline
	\finishhead
	cmd\_flash & device, firmware, enable\_insecure, log\_cb, progress\_cb, complete\_cb, error\_cb & Выполняет прошивку: проверка, отправка, перезагрузка. \\
	\hline
	on\_click\_flash & Нет аргументов & Обработчик кнопки «Flash» (чтение файла, разблокировка, запуск потока). \\
	\hline
	log & line & Добавляет строку в лог с временной меткой. \\
	\hline
	lock\_ui / unlock\_ui & force\_refresh & Блокирует элементы управления на время прошивки. \\
	\hline
	find\_device\_with\newline \_uid & cls, uid & Ищет устройство (класс) по UID (используется после перезагрузки). \\
	\hline
\end{xltabular}

\subsection{Пакет \texttt{widgets}}

\subsubsection{Модуль \texttt{keyboard\_widget.py}}

Класс \texttt{KeyboardWidget} – основной виджет для отрисовки клавиатуры по KLE-описанию. Поддерживает масштабирование, выделение клавиш, отображение нажатых/зажатых состояний, маскированных кодов (два уровня текста). Содержит внутренние классы \texttt{KeyWidget} и \texttt{EncoderWidget}.

Методы класса \texttt{KeyboardWidget} представлены в таблице \ref{keyboardwidget:table}.

\begin{xltabular}{\textwidth}{|p{3.5cm}|p{3.5cm}|>{\setlength{\baselineskip}{0.7\baselineskip}}X|}
	\caption{Методы класса \texttt{KeyboardWidget}\label{keyboardwidget:table}}\\
	\hline \centrow Название метода & \centrow Аргументы & \centrow Описание \\
	\hline
	\finishhead
	set\_keys & keys, encoders & Загружает описание клавиш и энкодеров для отображения. \\
	\hline
	update\_layout & Нет аргументов & Пересчитывает позиции всех виджетов в соответствии с текущим масштабом и выбранными опциями. \\
	\hline
	paintEvent & event & Рисует клавиатуру с учётом состояний (активная, нажатая, подсвеченная). \\
	\hline
	hit\_test & pos & Определяет, какая клавиша и её часть (маска или нет) находится в точке \texttt{pos}. \\
	\hline
	select\_next & Нет аргументов & Переключает выбор на следующую клавишу (по порядку). \\
	\hline
\end{xltabular}

\subsubsection{Модуль \texttt{key\_widget.py}}

Класс \texttt{KeyWidget} – компактный виджет для выбора клавиши. При клике открывает панель \texttt{KeycodeTabs}. Поддерживает фильтрацию кодов (только базовые для внутренней части маскированной клавиши).

Методы класса \texttt{KeyWidget} представлены в таблице \ref{kewwidget:table}.

\begin{xltabular}{\textwidth}{|p{3.5cm}|p{3.5cm}|>{\setlength{\baselineskip}{0.7\baselineskip}}X|}
	\caption{Методы класса \texttt{KeyWidget}\label{kewwidget:table}}\\
	\hline \centrow Название метода & \centrow Аргументы & \centrow Описание \\
	\hline
	\finishhead
	set\_keycode & kc & Устанавливает код клавиши и обновляет отображение. \\
	\hline
	set\_keycode\_\newline filter & keycode\_filter & Устанавливает функцию-фильтр для допустимых кодов. \\
	\hline
	on\_keycode\_\newline changed & keycode & Обрабатывает выбор кода из панели (учитывает маскирование). \\
	\hline
	mousePressEvent & ev & Открывает панель \texttt{KeycodeTabs} при клике. \\
	\hline
\end{xltabular}

\subsubsection{Модуль \texttt{square\_button.py}}

Класс \texttt{SquareButton} (наследник \texttt{QPushButton}) с квадратным sizeHint, масштабируемым относительно высоты шрифта. Поддерживает перенос текста (word wrap) через вложенный \texttt{QLabel}.

Методы класса \texttt{SquareButton} представлены в таблице \ref{square-button:table}.

\begin{xltabular}{\textwidth}{|p{3.5cm}|p{3.5cm}|>{\setlength{\baselineskip}{0.7\baselineskip}}X|}
	\caption{Методы класса \texttt{SquareButton}\label{square-button:table}}\\
	\hline \centrow Название метода & \centrow Аргументы & \centrow Описание \\
	\hline
	\finishhead
	setRelSize & ratio & Задаёт относительный размер кнопки относительно высоты шрифта. \\
	\hline
	setWordWrap & state & Включает/отключает перенос текста. \\
	\hline
	sizeHint & Нет аргументов & Возвращает квадратный размер кнопки. \\
	\hline
\end{xltabular}

\subsubsection{Модуль \texttt{flow\_layout.py}}

Класс \texttt{FlowLayout} – менеджер компоновки, который автоматически переносит элементы на следующую строку при нехватке места по горизонтали. Используется в панелях кнопок ключей и в редакторах макросов.

Методы класса \texttt{FlowLayout} представлены в таблице \ref{flow-layout:table}.

\begin{xltabular}{\textwidth}{|p{3.5cm}|p{3.5cm}|>{\setlength{\baselineskip}{0.7\baselineskip}}X|}
	\caption{Методы класса \texttt{FlowLayout}\label{flow-layout:table}}\\
	\hline \centrow Название метода & \centrow Аргументы & \centrow Описание \\
	\hline
	\finishhead
	addItem & item & Добавляет элемент в компоновку. \\
	\hline
	heightForWidth & width & Вычисляет высоту, необходимую для размещения всех элементов при заданной ширине. \\
	\hline
	doLayout & rect, testOnly & Выполняет фактическое размещение элементов. \\
	\hline
\end{xltabular}

\subsubsection{Модуль \texttt{tab\_widget\_keycodes.py}}

Класс \texttt{TabWidgetWithKeycodes} – наследник \texttt{QTabWidget}, который закрывает панель \texttt{KeycodeTabs} при переключении вкладки или клике вне её.

Методы класса  представлены в таблице \texttt{TabWidgetWithKeycodes\ref{tab-keycodes:table}}.

\begin{xltabular}{\textwidth}{|p{3.5cm}|p{3.5cm}|>{\setlength{\baselineskip}{0.7\baselineskip}}X|}
	\caption{Методы класса \texttt{TabWidgetWithKeycodes\label{tab-keycodes:table}}}\\
	\hline \centrow Название метода & \centrow Аргументы & \centrow Описание \\
	\hline
	\finishhead
	mouseRelease\newline Event & ev & Закрывает панель \texttt{KeycodeTabs}. \\
	\hline
	on\_changed & index & Закрывает панель \texttt{KeycodeTabs} при смене вкладки. \\
	\hline
\end{xltabular}

\subsubsection{Вспомогательные виджеты}

Модули \texttt{checkbox\_no\_padding.py} (убирает отступы у \texttt{QCheckBox}), \texttt{clickable\_label.py} (\texttt{QLabel} с сигналом \texttt{clicked}), \texttt{display\_keyboard.py} (статическое отображение клавиатуры для выбора кодов), \texttt{editor\_container.py} (контейнер для вкладок, передающий клик редактору). Каждый из них реализует простую специализированную функциональность, таблицы для них не приводятся ввиду тривиальности.

\subsection{Пакет \texttt{macro}}

\subsubsection{Модуль \texttt{macro\_action.py}}

Определяет классы действий макроса: \texttt{ActionText}, \texttt{ActionTap}, \texttt{ActionDown}, \texttt{ActionUp}, \texttt{ActionDelay}. Каждый содержит методы \texttt{serialize} (в бинарный вид), \texttt{save} (в список для JSON), \texttt{restore}.

Методы классов действий макросов представлены в таблице \ref{macro-action:table}.

\begin{xltabular}{\textwidth}{|p{3.5cm}|p{3.5cm}|>{\setlength{\baselineskip}{0.7\baselineskip}}X|}
	\caption{Методы классов действий макросов\label{macro-action:table}}\\
	\hline \centrow Название метода & \centrow Аргументы & \centrow Описание \\
	\hline
	\finishhead
	serialize & comm\_protocol & Преобразует действие в бинарный буфер для отправки на устройство. \\
	\hline
	save & Нет аргументов & Возвращает список, сериализуемый в JSON. \\
	\hline
	restore & act & Восстанавливает действие из JSON-данных. \\
	\hline
	\_\_eq\_\_ & other & Сравнивает два действия. \\
	\hline
\end{xltabular}

\subsubsection{Модуль \texttt{macro\_action\_ui.py}}

Содержит классы UI для каждого типа действия (\texttt{ActionTextUI}, \texttt{ActionTapUI}, …). Они создают соответствующие виджеты (поле ввода для текста, последовательность кнопок для \texttt{Tap/Down/Up}, спинбокс для задержки). Используются в \texttt{MacroLine}.

Классы UI действий макросов представлены в таблице \ref{macro-action-ui:table}.

\begin{xltabular}{\textwidth}{|p{3.5cm}|p{3.5cm}|>{\setlength{\baselineskip}{0.7\baselineskip}}X|}
	\caption{Классы UI действий макросов\label{macro-action-ui:table}}\\
	\hline \centrow Название класса & \centrow Методы & \centrow Описание \\
	\hline
	\finishhead
	ActionTextUI & \texttt{\_\_init\_\_, insert, remove, delete, on\_change} & Редактор текстовой строки. \\
	\hline
	ActionSequence\newline UI & \texttt{\_\_init\_\_, recreate\_sequence, on\_add, on\_change} & Базовый класс для последовательностей клавиш (Tap/Down/Up). \\
	\hline
	ActionDelayUI & \texttt{\_\_init\_\_, insert, remove, delete, on\_change} & Редактор задержки (SpinBox). \\
	\hline
\end{xltabular}

\subsubsection{Модуль \texttt{macro\_key.py}}

Классы для представления захваченных при записи нажатий: \texttt{KeyDown}, \texttt{KeyUp}, \texttt{KeyTap}, \texttt{KeyString}. Служат промежуточным звеном между рекордером и действиями.

Классы модуля \texttt{macro\_key.py} представлены в таблице \ref{macro-key:table}.

\begin{xltabular}{\textwidth}{|p{3.5cm}|p{3.5cm}|>{\setlength{\baselineskip}{0.7\baselineskip}}X|}
	\caption{Классы модуля \texttt{macro\_key.py}\label{macro-key:table}}\\
	\hline \centrow Название класса & \centrow Атрибуты & \centrow Описание \\
	\hline
	\finishhead
	KeyDown & keycode & Представляет нажатие клавиши. \\
	\hline
	KeyUp & keycode & Представляет отпускание клавиши. \\
	\hline
	KeyTap & keycode & Представляет нажатие и немедленное отпускание. \\
	\hline
	KeyString & string & Представляет последовательность печатных символов. \\
	\hline
\end{xltabular}

\subsubsection{Модуль \texttt{macro\_optimizer.py}}

Функции оптимизации макросов: \texttt{remove\_repeats} (удаление дублирующихся нажатий), \texttt{replace\_with\_tap} (замена последовательности Down+Up на Tap), \texttt{replace\_with\_string} (замена нескольких Tap печатных символов на один \texttt{KeyString}).

Функции модуля \texttt{macro\_optimizer.py} представлены в таблице \ref{macro-optimizer:table}.

\begin{xltabular}{\textwidth}{|p{3.5cm}|p{3.5cm}|>{\setlength{\baselineskip}{0.7\baselineskip}}X|}
	\caption{Функции модуля \texttt{macro\_optimizer.py}\label{macro-optimizer:table}}\\
	\hline \centrow Название функции & \centrow Аргументы & \centrow Описание \\
	\hline
	\finishhead
	remove\_repeats & sequence & Удаляет подряд идущие одинаковые нажатия или отпускания. \\
	\hline
	replace\_with\_tap & sequence & Заменяет пары Down+Up на Tap. \\
	\hline
	replace\_with\_\newline string & sequence & Группирует подряд идущие печатные Tap в KeyString. \\
	\hline
	macro\_optimize & sequence & Применяет все три оптимизации последовательно. \\
	\hline
\end{xltabular}

\subsubsection{Модуль \texttt{macro\_recorder\_linux.py}}

Класс \texttt{LinuxRecorder} (только для Linux) запускает привилегированный процесс через \texttt{pkexec}, который перехватывает глобальные нажатия клавиш с помощью библиотеки \texttt{keyboard}. Полученные события преобразуются в \texttt{KeyDown/KeyUp} и передаются сигналом \texttt{keystroke}.

Методы класса \texttt{LinuxRecorder} представлены в таблице \ref{linuxrecorder:table}.

\begin{xltabular}{\textwidth}{|p{3.5cm}|p{3.5cm}|>{\setlength{\baselineskip}{0.7\baselineskip}}X|}
	\caption{Методы класса \texttt{LinuxRecorder}\label{linuxrecorder:table}}\\
	\hline \centrow Название метода & \centrow Аргументы & \centrow Описание \\
	\hline
	\finishhead
	start & Нет аргументов & Запускает процесс записи (показывает окно «Stop recording»). \\
	\hline
	stop & Нет аргументов & Останавливает запись и завершает процесс. \\
	\hline
	on\_output & Нет аргументов & Читает строки из stdout процесса, парсит события и генерирует сигнал \texttt{keystroke}. \\
	\hline
\end{xltabular}

\subsubsection{Модуль \texttt{macro\_tab.py}}

Класс \texttt{MacroTab} – вкладка редактора одного макроса. Содержит кнопки «Record macro», «Add action», «Tap Enter», «Open Text Editor». Управляет списком строк (\texttt{MacroLine}), сериализует/десериализует макрос в JSON для текстового редактора.

Методы класса \texttt{MacroTab} представлены в таблице \ref{macro-tab:table}.

\begin{xltabular}{\textwidth}{|p{3.5cm}|p{3.5cm}|>{\setlength{\baselineskip}{0.7\baselineskip}}X|}
	\caption{Методы класса \texttt{MacroTab}\label{macro-tab:table}}\\
	\hline \centrow Название метода & \centrow Аргументы & \centrow Описание \\
	\hline
	\finishhead
	add\_action & act & Добавляет новое действие в конец макроса. \\
	\hline
	clear & Нет аргументов & Удаляет все действия из макроса. \\
	\hline
	on\_move & obj, offset & Перемещает строку действия вверх/вниз. \\
	\hline
	on\_text\_window & Нет аргументов & Открывает текстовый редактор для JSON-представления макроса. \\
	\hline
	actions & Нет аргументов & Возвращает список действий макроса. \\
	\hline
\end{xltabular}

\subsubsection{Модуль \texttt{macro\_line.py}}

Класс \texttt{MacroLine} – строка редактора макроса: содержит кнопки перемещения (▲/▼), выпадающий список типа действия, кнопку удаления (×) и виджет самого действия.

Методы класса \texttt{MacroLine} представлены в таблице \ref{macro-line:table}.

\begin{xltabular}{\textwidth}{|p{3.5cm}|p{3.5cm}|>{\setlength{\baselineskip}{0.7\baselineskip}}X|}
	\caption{Методы класса \texttt{MacroLine}\label{macro-line:table}}\\
	\hline \centrow Название метода & \centrow Аргументы & \centrow Описание \\
	\hline
	\finishhead
	insert & row & Вставляет строку в компоновку редактора. \\
	\hline
	remove & Нет аргументов & Удаляет строку из компоновки. \\
	\hline
	on\_change\_type & Нет аргументов & Заменяет тип действия при выборе в комбобоксе. \\
	\hline
	on\_remove\_\newline clicked & Нет аргументов & Удаляет строку. \\
	\hline
	on\_move\_up / on\_move\_down & Нет аргументов & Перемещает строку. \\
	\hline
\end{xltabular}

\subsection{Модуль \texttt{keycode\_tabs.py}}

Класс \texttt{KeycodeTabs} – панель с вкладками для выбора кодов клавиш. Содержит \texttt{FilteredKeycodeTabs} (с фильтром всех кодов или только базовых). При клике на клавишу генерирует сигнал \texttt{keycode\_changed}. Поддерживает режим «tray» – плавающее окно, которое появляется при выборе \texttt{KeyWidget}.

Методы класса \texttt{KeycodeTabs} представлены в таблице \ref{keycode-tabs:table}.

\begin{xltabular}{\textwidth}{|p{3.5cm}|p{3.5cm}|>{\setlength{\baselineskip}{0.7\baselineskip}}X|}
	\caption{Методы класса \texttt{KeycodeTabs}\label{keycode-tabs:table}}\\
	\hline \centrow Название метода & \centrow Аргументы & \centrow Описание \\
	\hline
	\finishhead
	recreate\_keycode\newline \_buttons & Нет аргументов & Перестраивает все кнопки клавиш в соответствии с текущим фильтром. \\
	\hline
	set\_keycode\_\newline filter & keycode\_filter & Устанавливает фильтр (все коды / только базовые). \\
	\hline
	open\_tray & target, keycode\_filter & Открывает плавающее окно выбора клавиш для целевого виджета. \\
	\hline
	close\_tray & Нет аргументов & Закрывает плавающее окно. \\
	\hline
\end{xltabular}

\subsection{Модуль \texttt{themes.py}}

Класс \texttt{Theme} предоставляет набор цветовых палитр (Light, Dark, Arc, Nord, Dracula, Catppuccin и др.). Метод \texttt{set\_theme} применяет палитру ко всему приложению через \texttt{QApplication.setPalette}. Настройка сохраняется в \texttt{QSettings}.

Методы класса \texttt{Theme} представлены в таблице \ref{themes:table}.

\begin{xltabular}{\textwidth}{|p{3.5cm}|p{3.5cm}|>{\setlength{\baselineskip}{0.7\baselineskip}}X|}
	\caption{Методы класса \texttt{Theme}\label{themes:table}}\\
	\hline \centrow Название метода & \centrow Аргументы & \centrow Описание \\
	\hline
	\finishhead
	set\_theme & theme & Применяет цветовую палитру, соответствующую имени \texttt{theme}, ко всему приложению. \\
	\hline
	get\_theme & Нет аргументов & Возвращает имя текущей темы. \\
	\hline
	mask\_light\_factor & Нет аргументов & Возвращает коэффициент осветления для области маски (внутренней части маскированной клавиши). \\
	\hline
\end{xltabular}

\subsection{Модуль \texttt{security\_dialog.py}}

Класс \texttt{Unlocker} – модальный диалог, требующий от пользователя одновременного нажатия нескольких физических клавиш для перевода клавиатуры в разблокированное состояние.

Методы класса \texttt{Unlocker} представлены в таблице \ref{security:table}.

\begin{xltabular}{\textwidth}{|p{3.5cm}|p{3.5cm}|>{\setlength{\baselineskip}{0.7\baselineskip}}X|}
	\caption{Методы класса \texttt{Unlocker}\label{security:table}}\\
	\hline \centrow Название метода & \centrow Аргументы & \centrow Описание \\
	\hline
	\finishhead
	unlock & keyboard & Статический метод: отображает диалог, ожидает завершения, возвращает успех. \\
	\hline
	perform\_unlock & Нет аргументов & Отправляет команду начала разблокировки, запускает таймер опроса. \\
	\hline
	unlock\_poller & Нет аргументов & Периодически опрашивает статус, обновляет прогресс-бар, закрывает диалог при успехе. \\
	\hline
	update\_reference & Нет аргументов & Подсвечивает на миниатюре клавиши, которые необходимо удерживать. \\
	\hline
\end{xltabular}

\subsection{Модуль \texttt{utils.py}}

Содержит вспомогательные функции:
\begin{itemize}
	\item \texttt{hid\_send} – отправка HID-отчёта (дублирует одноимённую из \texttt{hid\_comm.py}, но используется в разных контекстах);
	\item \texttt{find\_custom\_devices} – поиск устройств по magic-строкам, legacy-стеку, sideload;
	\item \texttt{chunks} – разбиение буфера на части;
	\item \texttt{pad\_for\_bootloader} – дополнение сообщения до 64 байт нулями;
	\item \texttt{init\_logger} – настройка ротируемого логгера;
	\item \texttt{make\_scrollable} – оборачивание компоновки в \texttt{QScrollArea};
	\item \texttt{KeycodeDisplay} – класс для отображения текста на клавишах с учётом выбранной раскладки и маскированных кодов.
\end{itemize}

Функции и классы модуля \texttt{utils.py} представлены в таблице \ref{utils:table}.

\begin{xltabular}{\textwidth}{|p{3.5cm}|p{3.5cm}|>{\setlength{\baselineskip}{0.7\baselineskip}}X|}
	\caption{Функции и классы модуля \texttt{utils.py}\label{utils:table}}\\
	\hline \centrow Название & \centrow Аргументы & \centrow Описание \\
	\hline
	\finishhead
	find\_custom\_\newline devices & legacy\_stack, sideload\_vid, sideload\_pid, quiet & Возвращает список объектов устройств. \\
	\hline
	chunks & data, sz & Генератор, разбивающий \texttt{data} на куски размера \texttt{sz}. \\
	\hline
	pad\_for\_\newline bootloader & msg & Дополняет сообщение нулями до 64 байт. \\
	\hline
	init\_logger & Нет аргументов & Настраивает логгирование в файл. \\
	\hline
	make\_scrollable & layout & Оборачивает компоновку в \texttt{QScrollArea}. \\
	\hline
	KeycodeDisplay\newline .get\_label & code & Возвращает текст для отображения на клавише. \\
	\hline
	KeycodeDisplay\newline .set\_keymap\_\newline override & override & Устанавливает активную раскладку отображения. \\
	\hline
\end{xltabular}

\subsection{Модуль \texttt{main\_window.py}}

Главное окно приложения (\texttt{MainWindow}) содержит комбобокс выбора устройства, панель вкладок с редакторами, меню (Файл, Безопасность, Тема, Справка).

Методы класса \texttt{MainWindow} представлены в таблице \ref{tab:mainwindow}.

\begin{xltabular}{\textwidth}{|p{3.5cm}|p{3.5cm}|>{\setlength{\baselineskip}{0.7\baselineskip}}X|}
	\caption{Методы класса \texttt{MainWindow}\label{tab:mainwindow}}\\
	\hline \centrow Название метода & \centrow Аргументы & \centrow Описание \\
	\hline
	\endfirsthead
	\caption*{Продолжение таблицы \ref{tab:mainwindow}}\\
	\hline \centrow Название метода & \centrow Аргументы & \centrow Описание \\
	\hline
	\finishhead
	on\_devices\_\newline updated & devices, hard\_refresh & Обновляет комбобокс устройств, скрывает/показывает сообщение об отсутствии устройств. \\
	\hline
	on\_device\_\newline selected & Нет аргументов & Переключает текущее устройство, перестраивает все редакторы. \\
	\hline
	rebuild & Нет аргументов & Перезагружает все редакторы для выбранного устройства. \\
	\hline
	refresh\_tabs & Нет аргументов & Обновляет набор вкладок согласно поддерживаемым возможностям клавиатуры. \\
	\hline
	change\_\newline keyboard\_\newline layout & index & Изменяет языковую раскладку отображения клавиш. \\
	\hline
	set\_theme & theme & Применяет выбранную цветовую тему. \\
	\hline
	unlock\_keyboard / lock\_keyboard & Нет аргументов & Вызывает диалог разблокировки или отправляет команду блокировки. \\
	\hline
	on\_layout\_load / on\_layout\_save & Нет аргументов & Загружает/сохраняет файлы \texttt{.vil} (раскладка, макросы, настройки). \\
	\hline
\end{xltabular}

\subsection{Модуль \texttt{text\_editor.py}}

Класс \texttt{TextboxWindow} – диалог с многострочным текстовым полем, предназначенный для редактирования JSON-представления макросов. Поддерживает импорт/экспорт файлов, копирование/вставку.

Методы класса \texttt{TextboxWindow} представлены в таблице \ref{TextboxWindow:table}.

\begin{xltabular}{\textwidth}{|p{3.5cm}|p{3.5cm}|>{\setlength{\baselineskip}{0.7\baselineskip}}X|}
	\caption{Методы класса \texttt{TextboxWindow}\label{TextboxWindow:table}}\\
	\hline \centrow Название метода & \centrow Аргументы & \centrow Описание \\
	\hline
	\finishhead
	getText & Нет аргументов & Возвращает содержимое текстового поля. \\
	\hline
	on\_apply & Нет аргументов & Закрывает диалог с кодом \texttt{Accept}. \\
	\hline
	on\_copy / on\_paste & Нет аргументов & Копирует/вставляет текст. \\
	\hline
	on\_export / on\_import & Нет аргументов & Сохраняет/загружает текст в файл. \\
	\hline
\end{xltabular}

\subsection{Модули \texttt{about\_device.py}, \texttt{custom\_code.py}, \texttt{custom\_code\_dialog.py}}

\begin{itemize}
	\item \texttt{about\_device.py} – диалог «О устройстве», отображающий информацию о VID/PID, версиях протоколов, количестве макросов, tap dance, комбинаций и поддержке фич.
	\item \texttt{custom\_code.py} – функции для создания сложных кодов клавиш (например, \texttt{LSFT(KC\_A)}) с использованием синтаксического анализатора \texttt{simpleeval}. Определяет функции-обёртки: \texttt{LCTL}, \texttt{LSFT}, \texttt{MT}, \texttt{LT}, \texttt{TD} и т.п.
	\item \texttt{custom\_code\_dialog.py} – диалог для ручного ввода произвольного выражения (например, \texttt{LSFT(KC\_B)}). Отображает вычисленное значение и проверяет корректность.
\end{itemize}

Классы и функции модулей \texttt{about\_device.py}, \texttt{custom\_code.py}, \texttt{custom\_code\_dialog.py} представлены в таблице \ref{about-device:table}.

\begin{xltabular}{\textwidth}{|p{3.5cm}|p{3.5cm}|>{\setlength{\baselineskip}{0.7\baselineskip}}X|}
	\caption{Классы и функции модулей \texttt{about\_device.py}, \texttt{custom\_code.py}, \texttt{custom\_code\_dialog.py}\label{about-device:table}}\\
	\hline \centrow Название & \centrow Аргументы & \centrow Описание \\
	\hline
	\finishhead
	AboutDevice & device & Конструктор, заполняет диалог информацией об устройстве. \\
	\hline
	CustomKeyCode\newline .decode & s & Разбирает строку-выражение (например, \texttt{"LSFT(KC\_A)"}) в числовой код клавиши. \\
	\hline
	CustomCode\newline Dialog & initial & Диалог для ручного ввода выражения, отображает вычисленное значение. \\
	\hline
\end{xltabular}

\subsection{Модуль \texttt{layout\_presets.py} и пакет \texttt{layouts}}

Модуль \texttt{layout\_presets.py} импортирует словари переопределения надписей из файлов пакета \texttt{layouts} (например, \texttt{german.py}, \texttt{french.py}, \texttt{colemak.py} и др.). Экспортирует список \texttt{LAYOUT\_PRESETS} – кортежи (имя раскладки, словарь отображений QMK ID → строка). Выбранная раскладка применяется через \texttt{KeycodeDisplay.set\_keymap\_override}. Таблицы не требуются, так как модуль не содержит функций или классов – только данные.

\subsection{Тестирование программной системы}

В ходе системного тестирования программного продукта, предназначенного для настройки сплит-клавиатур, применялись разные методы и инструменты. Результаты проведенного тестирования отражены в данном разделе посредством демонстраций экранных изображений выполненных операций.

Функция: загрузка файла конфигурации.
Условия: пользователь должен перейти на вкладку "<File">, выбрать файл формата .vial (желательно ранее созданный профиль) с диска своего компьютера и загрузить его в программу.
Входные данные: выбранный пользователем файл конфигурации (.vial), содержащий ранее настроенную раскладку клавиш, слои и макросы.
Результат: программа обработает указанный файл и восстановит предыдущую конфигурацию клавиатуры, включающую всю сохранённую информацию о расположении клавиш, слоях, макросах и анимации подсветки. После загрузки изменения сразу вступают в силу, и клавиатура готова к работе с новыми параметрами (рисунки \ref{test1_1:image} -- \ref{test1_2:image}). 

\begin{figure}[H]
	\centering
	\includegraphics[width=0.75\linewidth]{Загрузка файла конфигурации.png}
	\caption{Загрузка файла конфигурации}
	\label{test1_1:image}
\end{figure}

\begin{figure}[H]
	\centering
	\includegraphics[width=0.75\linewidth]{Результат загрузки файла конфигурации.png}
	\caption{Результат загрузки файла конфигурации}
	\label{test1_2:image}
\end{figure}

Функция: настройка раскладки клавиатуры.
Условия: пользователь должен перейти на вкладку "<Keyboard Layout">, выбрать нужную схему расположения клавиш (например, стандартный ANSI или альтернативные варианты вроде Colemak или Dvorak), задать слой (Layer), к которому относятся кнопки, и определить назначение каждой клавиши индивидуально.
Входные данные: предпочтения пользователя относительно схемы размещения клавиш, выбранного слоя и индивидуальных настроек отдельных кнопок.
Результат: приложение сохраняет внесённые изменения в память клавиатуры, обновляя её текущую раскладку. После сохранения новая схема становится активной немедленно, и клавиатура перестаёт реагировать на предыдущие установки. Пользователь получает возможность полноценно пользоваться выбранной раскладкой и активированными функциями, такими как смена слоя или запуск макросов (рисунки \ref{test2_1:image} -- \ref{test2_4:image}). 

\begin{figure}[H]
	\centering
	\includegraphics[width=0.75\linewidth]{Результат настройки раскладки клавиатуры (нулевой слой).png}
	\caption{Результат настройки раскладки клавиатуры (нулевой слой)}
	\label{test2_1:image}
\end{figure}

\begin{figure}[H]
	\centering
	\includegraphics[width=0.75\linewidth]{Результат настройки раскладки клавиатуры (первый слой).png}
	\caption{Результат настройки раскладки клавиатуры (первый слой)}
	\label{test2_2:image}
\end{figure}

\begin{figure}[H]
	\centering
	\includegraphics[width=0.75\linewidth]{Результат настройки раскладки клавиатуры (второй слой).png}
	\caption{Результат настройки раскладки клавиатуры (второй слой)}
	\label{test2_3:image}
\end{figure}

\begin{figure}[H]
	\centering
	\includegraphics[width=0.75\linewidth]{Результат настройки раскладки клавиатуры (третий слой).png}
	\caption{Результат настройки раскладки клавиатуры (третий слой)}
	\label{test2_4:image}
\end{figure}

Функция: защита конфигурации клавиатуры.
Условия: пользователь должен перейти на вкладку "<Security">, нажать Unlock, зажать определённые клавиши, которые будут подсвечены в окне, подождать, пока клавиатура не будет разблокирована для дальнейших взаимодействий.
Входные данные: нажатие кнопки Unlock, зажатые клавиши
Результат: клавиатура разблокирована для дальнейших взаимодействий (рисунки \ref{test3_1:image} -- \ref{test3_3:image}).

\begin{figure}[H]
	\centering
	\includegraphics[width=0.75\linewidth]{защита конфигурации клавиатуры (клавиатура заблокирована).png}
	\caption{защита конфигурации клавиатуры (клавиатура заблокирована)}
	\label{test3_1:image}
\end{figure}

\begin{figure}[H]
	\centering
	\includegraphics[width=0.75\linewidth]{Процесс разблокировки клавиатуры.png}
	\caption{Процесс разблокировки клавиатуры}
	\label{test3_2:image}
\end{figure}

\begin{figure}[H]
	\centering
	\includegraphics[width=0.75\linewidth]{Клавиатура разблокирована.png}
	\caption{Клавиатура разблокирована}
	\label{test3_3:image}
\end{figure}

Функция: выбор цветовой темы.
Условия: пользователь должен перейти на вкладку "<Theme">, выбрать понравившуюся цветовую палитру из предложенных вариантов или настроить собственную комбинацию цветов.
Входные данные: предпочтительная цветовая схема или индивидуальный подбор оттенков подсветки клавиатуры.
Результат: после выбора или задания собственных предпочтений, клавиатура моментально принимает новую расцветку подсветки, создавая привлекательное световое оформление, соответствующее выбору пользователя. Новая тема сохраняется в памяти устройства и применяется постоянно до следующего изменения (рисунки \ref{test4_1:image} -- \ref{test4_2:image}).

\begin{figure}[H]
	\centering
	\includegraphics[width=0.75\linewidth]{Выбор цветовой темы.png}
	\caption{Выбор цветовой темы}
	\label{test4_1:image}
\end{figure}

\begin{figure}[H]
	\centering
	\includegraphics[width=0.75\linewidth]{Результат смены цветовой темы.png}
	\caption{Результат смены цветовой темы}
	\label{test4_2:image}
\end{figure}

Функция: создание собственной карты клавиш (Keymap).
Условия: пользователь должен перейти на вкладку "<Keymap">, выбрать нужный слой (Layer) и назначить каждому ключу требуемое значение (символ, команду или макрос).
Входные данные: уникальные значения, указанные пользователем для каждого ключа на каждом слое.
Результат: после завершения процесса программирования ключей и нажатия кнопки "<Save">, созданная карта клавиш сохраняется в памяти клавиатуры. Она вступает в силу незамедлительно, и все последующие нажатия клавиш соответствуют новым назначенным командам. Таким образом, пользователь получает возможность использовать свою уникальную карту клавиш, соответствующую индивидуальным потребностям и привычкам (рисунки \ref{test5_1:image} -- \ref{test5_2:image}).

\begin{figure}[H]
	\centering
	\includegraphics[width=0.75\linewidth]{Настройка раскладки клавиатуры.png}
	\caption{Настройка раскладки клавиатуры}
	\label{test5_1:image}
\end{figure}

\begin{figure}[H]
	\centering
	\includegraphics[width=0.75\linewidth]{Результат смены цветовой темы.png}
	\caption{Результат смены цветовой темы}
	\label{test5_2:image}
\end{figure}

Функция: создание и управление макросами.
Условия: пользователь должен перейти на вкладку "<Macros">, выбрать необходимую категорию макросов (текстовые команды, мультимедийные функции, игровые комбинации и др.) и задать конкретные последовательности действий, которые будут выполняться при нажатии определённых клавиш.
Входные данные: комбинация клавиш и последовательность выполняемых действий, указанных пользователем.
Результат: после создания и сохранения макроса, клавиатура запоминает установленную последовательность. Теперь каждое последующее нажатие указанной комбинации вызывает немедленное исполнение записанных команд. Макросы значительно ускоряют повторяющиеся операции и повышают эффективность работы пользователя (рисунки \ref{test6_1:image} -- \ref{test6_2:image}). 

\begin{figure}[H]
	\centering
	\includegraphics[width=0.75\linewidth]{Создание макроса.png}
	\caption{Создание макроса}
	\label{test6_1:image}
\end{figure}

\begin{figure}[H]
	\centering
	\includegraphics[width=0.75\linewidth]{Результат создания макроса.png}
	\caption{Результат создания макроса}
	\label{test6_2:image}
\end{figure}

Функция: настройка подсветки клавиатуры.
Условия: пользователь должен перейти на вкладку "<Lighting">, выбрать одну из предустановленных схем освещения или вручную задать интенсивность, цвет и эффекты подсветки клавиш.
Входные данные: предпочтительные параметры подсветки, установленные пользователем.
Результат: после внесения изменений в настройки подсветки, клавиатура мгновенно примет указанную пользователем схему. Светодиоды будут светиться нужным цветом и яркостью, подчеркивая дизайн и улучшая эстетику устройства. Новые параметры сохраняются в памяти клавиатуры и применяются автоматически каждый раз при включении питания (рисунки \ref{test7_1:image} -- \ref{test7_3:image}).

\begin{figure}[H]
	\centering
	\includegraphics[width=0.75\linewidth]{Настройка цвета подсветки клавиатуры.png}
	\caption{Настройка цвета подсветки клавиатуры}
	\label{test7_1:image}
\end{figure}

\begin{figure}[H]
	\centering
	\includegraphics[width=0.75\linewidth]{Настройка режима подсветки клавиатуры.png}
	\caption{Настройка режима подсветки клавиатуры}
	\label{test7_2:image}
\end{figure}

\begin{figure}[H]
	\centering
	\includegraphics[width=0.75\linewidth]{Результат настройки подсветки.png}
	\caption{Результат настройки подсветки}
	\label{test7_3:image}
\end{figure}
 
Функция: настройка поведения тап-дансов (Tap Dance).
Условия: пользователь должен перейти на вкладку "<Tap Dance">, выбрать интересующую клавишу или сочетание клавиш и задать желаемую реакцию на количество быстрых последовательных нажатий (одинарное, двойное, тройное касание и т.п.).
Входные данные: ключи, выбранные пользователем, и задаваемое ими поведение в зависимости от количества быстрых нажатий.
Результат: после сохранения настроек, клавиатура начнёт распознавать указанное число быстрых нажатий одной клавиши и выполнять соответствующие действия. Например, однократное нажатие клавиши может привести к одному действию, двукратное -- к другому, и так далее. Эти реакции облегчают быстрый доступ к часто используемым действиям и улучшают производительность пользователя (рисунки \ref{test8_1:image} -- \ref{test8_2:image}). 

\begin{figure}[H]
	\centering
	\includegraphics[width=0.75\linewidth]{Настройка поведения тап-дансов (Tap Dance).png}
	\caption{Настройка поведения тап-дансов (Tap Dance)}
	\label{test8_1:image}
\end{figure}

\begin{figure}[H]
	\centering
	\includegraphics[width=0.75\linewidth]{Результат настройки поведения тап-дансов (Tap Dance).png}
	\caption{Результат настройки поведения тап-дансов (Tap Dance)}
	\label{test8_2:image}
\end{figure}

Функция: создание комбинаций клавиш (Combos).
Условия: пользователь должен перейти на вкладку "<Combos">, выбрать нужные сочетания клавиш и задать действия, которые будут выполнены при одновременном нажатии указанных клавиш.
Входные данные: комбинации клавиш, заданные пользователем, и связанные с ними команды или функции.
Результат: после настройки и сохранения комбинации, клавиатура начнет интерпретировать одновременно нажатые клавиши как единую команду. Например, одновременное нажатие двух клавиш может вызывать определенную функцию операционной системы или специализированное приложение. Комбинации помогают повысить скорость работы и упростить взаимодействие с компьютером (рисунки \ref{test9_1:image} -- \ref{test9_2:image}).

\begin{figure}[H]
	\centering
	\includegraphics[width=0.75\linewidth]{Создание комбинаций клавиш (Combos).png}
	\caption{Создание комбинаций клавиш (Combos)}
	\label{test9_1:image}
\end{figure}

\begin{figure}[H]
	\centering
	\includegraphics[width=0.75\linewidth]{Результат настройки раскладки клавиатуры2.png}
	\caption{Результат настройки раскладки клавиатуры}
	\label{test9_2:image}
\end{figure}

Функция: работа с несколькими слоями клавиатуры (Layers).
Условия: пользователь должен перейти на вкладку "<Layers">, выбрать основной слой (обычно обозначается Layer 0) и дополнительно добавить или удалить вспомогательные слои (Layer 1, Layer 2 и т.д.). Для каждого слоя можно отдельно задать своё расположение клавиш и функциональность.
Входные данные: необходимые пользователю слои и их настройки.
Результат: после сохранения новых слоёв клавиатура сможет переключаться между разными наборами символов и функциональных клавиш. Переключение между слоями осуществляется путём нажатия специальных клавиш-модификаторов ("<Fn">), позволяющих временно или длительно переходить на другой слой. Такая гибкая организация улучшает адаптацию клавиатуры под различные сценарии использования (рисунки \ref{test10_1:image} -- \ref{test10_4:image}).

\begin{figure}[H]
	\centering
	\includegraphics[width=0.75\linewidth]{Работа с несколькими слоями клавиатуры (Layers) (нулевой слой).png}
	\caption{Работа с несколькими слоями клавиатуры (Layers) (нулевой слой)}
	\label{test10_1:image}
\end{figure}

\begin{figure}[H]
	\centering
	\includegraphics[width=0.75\linewidth]{Работа с несколькими слоями клавиатуры (Layers) (первый слой).png}
	\caption{Работа с несколькими слоями клавиатуры (Layers) (первый слой)}
	\label{test10_2:image}
\end{figure}

\begin{figure}[H]
	\centering
	\includegraphics[width=0.75\linewidth]{Работа с несколькими слоями клавиатуры (Layers) (второй слой).png}
	\caption{Работа с несколькими слоями клавиатуры (Layers) (второй слой)}
	\label{test10_3:image}
\end{figure}

\begin{figure}[H]
	\centering
	\includegraphics[width=0.75\linewidth]{Работа с несколькими слоями клавиатуры (Layers) (третий слой).png}
	\caption{Работа с несколькими слоями клавиатуры (Layers) (третий слой)}
	\label{test10_4:image}
\end{figure}

Функция: тестирование работоспособности клавиатуры (Matrox Tester).
Условия: пользователь должен перейти на вкладку "<Tester">, подключить клавиатуру и инициировать проверку состояния контактов клавиш и матрицы клавиатуры.
Входные данные: подключённая клавиатура, выбранная в списке устройств, и начатая процедура проверки.
Результат: программа осуществляет сканирование всех клавиш, выявляет возможные проблемы с контактами, наличие зависших или неисправных элементов. Тестировщик выводит отчет о состоянии клавиатуры, отмечая исправные и проблемные зоны. Если обнаружены неисправности, пользователь получит рекомендации по устранению неполадок или замене дефектных компонентов (рисунки \ref{test11_1:image} -- \ref{test11_2:image}).

\begin{figure}[H]
	\centering
	\includegraphics[width=0.75\linewidth]{Тестирование работоспособности клавиатуры.png}
	\caption{Тестирование работоспособности клавиатуры}
	\label{test11_1:image}
\end{figure}

\begin{figure}[H]
	\centering
	\includegraphics[width=0.75\linewidth]{Результат настройки раскладки клавиатуры3.png}
	\caption{Результат настройки раскладки клавиатуры}
	\label{test11_2:image}
\end{figure}

Функция: переопределение значений клавиш (Key Overrides).
Условия: пользователь должен перейти на вкладку "<Key Overrides">, выбрать целевую клавишу и задать новое действие или символ, который будет заменять стандартное поведение клавиши.
Входные данные: название целевой клавиши и новое значение, присвоенное ей пользователем.
Результат: после сохранения настроек, клавиатура начинает воспринимать выбранную клавишу иначе. Каждое последующее нажатие теперь будет приводить к выполнению заранее назначенной команды или вводу конкретного символа. Переопределение клавиш полезно для адаптации клавиатуры под специфичные потребности пользователя или для улучшения производительности труда (рисунки \ref{test12_1:image} -- \ref{test12_2:image}).

\begin{figure}[H]
	\centering
	\includegraphics[width=0.75\linewidth]{Переопределение значений клавиш (Key Overrides).png}
	\caption{Переопределение значений клавиш (Key Overrides)}
	\label{test12_1:image}
\end{figure}

\begin{figure}[H]
	\centering
	\includegraphics[width=0.75\linewidth]{Результат переопределение значений клавиш (Key Overrides).png}
	\caption{Результат переопределение значений клавиш (Key Overrides)}
	\label{test12_2:image}
\end{figure}

Функция: установка альтернативного повтора нажатия клавиши (Alt Repeat Key).
Условия: пользователь должен перейти на вкладку "<Alt Repeat Key">, выбрать ключевую клавишу, которую необходимо заменить функцией автоповтора, и указать интервал задержки перед началом повторного срабатывания.
Входные данные: выделенная клавиша и выбранный промежуток времени для начала автоматического повторения.
Результат: после сохранения настроек, данная клавиша начнёт повторять своё действие автоматически, начиная через установленный пользователем временной интервал. Такой подход полезен для автоматизации некоторых рутинных задач, таких как перемещение курсора или пролистывание текста (рисунки \ref{test13_1:image} -- \ref{test13_2:image}).

\begin{figure}[H]
	\centering
	\includegraphics[width=0.75\linewidth]{Установка альтернативного повтора нажатия клавиши (Alt Repeat Key).png}
	\caption{Установка альтернативного повтора нажатия клавиши (Alt Repeat Key)}
	\label{test13_1:image}
\end{figure}

\begin{figure}[H]
	\centering
	\includegraphics[width=0.75\linewidth]{Результат установки альтернативного повтора нажатия клавиши (Alt Repeat Key).png}
	\caption{Результат установки альтернативного повтора нажатия клавиши (Alt Repeat Key)}
	\label{test13_2:image}
\end{figure}

Функция: настройка общих параметров QMK-прошивки (QMK Settings).
Условия: пользователь должен перейти на вкладку "<QMK Settings">, выбрать интересующие параметры, такие как тайминги, скорость движения мыши, чувствительность колёс прокрутки и другие глобальные настройки.
Входные данные: предпочитаемые пользователем параметры QMK-прошивки.
Результат: после подтверждения настроек, изменения вносятся в прошивку клавиатуры и становятся активными немедленно. Общие настройки влияют на общее поведение клавиатуры, улучшив её отзывчивость и удобство использования (рисунки \ref{test14_1:image} -- \ref{test14_2:image}).

\begin{figure}[H]
	\centering
	\includegraphics[width=0.75\linewidth]{Настройка общих параметров QMK-прошивки (QMK Settings).png}
	\caption{Настройка общих параметров QMK-прошивки (QMK Settings)}
	\label{test14_1:image}
\end{figure}

\begin{figure}[H]
	\centering
	\includegraphics[width=0.75\linewidth]{Результат настройки общих параметров QMK-прошивки (QMK Settings).png}
	\caption{Результат настройки общих параметров QMK-прошивки (QMK Settings)}
	\label{test14_2:image}
\end{figure}
